\documentclass[a4paper,11pt]{book}
%\documentclass[a4paper,twoside,11pt,titlepage]{book}
\usepackage{listings}
\usepackage[utf8]{inputenc}
\usepackage[spanish]{babel}

% \usepackage[style=list, number=none]{glossary} %
%\usepackage{titlesec}
%\usepackage{pailatino}

\decimalpoint
\usepackage{dcolumn}
\newcolumntype{.}{D{.}{\esperiod}{-1}}
\makeatletter
\addto\shorthandsspanish{\let\esperiod\es@period@code}
\makeatother


%\usepackage[chapter]{algorithm}
\RequirePackage{verbatim}
%\RequirePackage[Glenn]{fncychap}
\usepackage{fancyhdr}
\usepackage{graphicx}
\usepackage{afterpage}

\usepackage{longtable}

% \usepackage[pdfborder={000}]{hyperref} %referencia
\usepackage{hyperref}

% ********************************************************************
% Re-usable information
% ********************************************************************
\newcommand{\myTitle}{SIGMA\xspace}
\newcommand{\mySubTitle}{Sistema Inteligente de Gestión Monetaria Automatizada\xspace}
\newcommand{\myDegree}{GRADO EN INGENIERÍA INFORMÁTICA\xspace}
\newcommand{\myName}{Elena Ortega Contreras (alumno)\xspace}
\newcommand{\myProf}{José Manuel Soto Hidalgo (tutor1)\xspace}
\newcommand{\myOtherProf}{Nombre Apllido1 Apellido2 (tutor2)\xspace}
%\newcommand{\mySupervisor}{Put name here\xspace}
\newcommand{\myFaculty}{Escuela Técnica Superior de Ingenierías Informática y de
Telecomunicación\xspace}
\newcommand{\myFacultyShort}{E.T.S. de Ingenierías Informática y de
Telecomunicación\xspace}
\newcommand{\myDepartment}{Departamento de ...\xspace}
\newcommand{\myUni}{\protect{Universidad de Granada}\xspace}
\newcommand{\myLocation}{Granada\xspace}
\newcommand{\myTime}{\today\xspace}
\newcommand{\myVersion}{Version 0.1\xspace}


\hypersetup{
pdfauthor = {\myName (email (en) ugr (punto) es)},
pdftitle = {\myTitle},
pdfsubject = {},
pdfkeywords = {palabra_clave1, palabra_clave2, palabra_clave3, ...},
pdfcreator = {LaTeX con el paquete ....},
pdfproducer = {pdflatex}
}

%\hyphenation{}


%\usepackage{doxygen/doxygen}
%\usepackage{pdfpages}
\usepackage{url}
\usepackage{colortbl,longtable}
\usepackage[stable]{footmisc}
%\usepackage{index}

%\makeindex
%\usepackage[style=long, cols=2,border=plain,toc=true,number=none]{glossary}
% \makeglossary

% Definición de comandos que me son tiles:
%\renewcommand{\indexname}{Índice alfabético}
%\renewcommand{\glossaryname}{Glosario}

\pagestyle{fancy}
\fancyhf{}
\fancyhead[LO]{\leftmark}
\fancyhead[RE]{\rightmark}
\fancyhead[RO,LE]{\textbf{\thepage}}
\renewcommand{\chaptermark}[1]{\markboth{\textbf{#1}}{}}
\renewcommand{\sectionmark}[1]{\markright{\textbf{\thesection. #1}}}

\setlength{\headheight}{1.5\headheight}

\newcommand{\HRule}{\rule{\linewidth}{0.5mm}}
%Definimos los tipos teorema, ejemplo y definición podremos usar estos tipos
%simplemente poniendo \begin{teorema} \end{teorema} ...
\newtheorem{teorema}{Teorema}[chapter]
\newtheorem{ejemplo}{Ejemplo}[chapter]
\newtheorem{definicion}{Definición}[chapter]

\definecolor{gray97}{gray}{.97}
\definecolor{gray75}{gray}{.75}
\definecolor{gray45}{gray}{.45}
\definecolor{gray30}{gray}{.94}

\lstset{ frame=Ltb,
     framerule=0.5pt,
     aboveskip=0.5cm,
     framextopmargin=3pt,
     framexbottommargin=3pt,
     framexleftmargin=0.1cm,
     framesep=0pt,
     rulesep=.4pt,
     backgroundcolor=\color{gray97},
     rulesepcolor=\color{black},
     %
     stringstyle=\ttfamily,
     showstringspaces = false,
     basicstyle=\scriptsize\ttfamily,
     commentstyle=\color{gray45},
     keywordstyle=\bfseries,
     %
     numbers=left,
     numbersep=6pt,
     numberstyle=\tiny,
     numberfirstline = false,
     breaklines=true,
   }
 
% minimizar fragmentado de listados
\lstnewenvironment{listing}[1][]
   {\lstset{#1}\pagebreak[0]}{\pagebreak[0]}

\lstdefinestyle{CodigoC}
   {
	basicstyle=\scriptsize,
	frame=single,
	language=C,
	numbers=left
   }
\lstdefinestyle{CodigoC++}
   {
	basicstyle=\small,
	frame=single,
	backgroundcolor=\color{gray30},
	language=C++,
	numbers=left
   }

 
\lstdefinestyle{Consola}
   {basicstyle=\scriptsize\bf\ttfamily,
    backgroundcolor=\color{gray30},
    frame=single,
    numbers=none
   }


\newcommand{\bigrule}{\titlerule[0.5mm]}


%Para conseguir que en las páginas en blanco no ponga cabecerass
\makeatletter
\def\clearpage{%
  \ifvmode
    \ifnum \@dbltopnum =\m@ne
      \ifdim \pagetotal <\topskip
        \hbox{}
      \fi
    \fi
  \fi
  \newpage
  \thispagestyle{empty}
  \write\m@ne{}
  \vbox{}
  \penalty -\@Mi
}
\makeatother

\usepackage{pdfpages}
\begin{document}
\begin{titlepage}
    \newlength{\centeroffset}
    \setlength{\centeroffset}{-0.5\oddsidemargin}
    \addtolength{\centeroffset}{0.5\evensidemargin}
    \thispagestyle{empty}
    
    \noindent\hspace*{\centeroffset}
    \begin{minipage}{\textwidth}
    
    \centering
    \includegraphics[width=0.9\textwidth]{logos/logo_ugr.jpg}\\[1.4cm]
    
    \textsc{ \Large TRABAJO FIN DE GRADO\\[0.2cm]}
    \textsc{ GRADO EN INGENIERIA INFORMATICA}\\[1cm]
    
    {\Huge\bfseries SIGMA \\}
    \noindent\rule[-1ex]{\textwidth}{3pt}\\[3.5ex]
    {\large\bfseries Sistema Inteligente de Gestión Monetaria Automatizada }
    \end{minipage}
    
    \vspace{2.5cm}
    \noindent\hspace*{\centeroffset}
    \begin{minipage}{\textwidth}
    \centering
    
    \textbf{Autor}\\ {Elena Ortega Contreras}\\[2.5ex]
    \textbf{Director}\\ {José Manuel Soto García}\\[2cm]
    \includegraphics[width=0.3\textwidth]{logos/etsiit_logo.png}\\[0.1cm]
    \textsc{Escuela Técnica Superior de Ingenierías Informática y de Telecomunicación}\\
    \textsc{---}\\
    Granada, Junio de 2024
    \end{minipage}
    \end{titlepage}
\thispagestyle{empty}

\begin{center}
{\large\bfseries Título \\ Subtítulo }\\
\end{center}
\begin{center}
Nombre Del Estudiante\\
\end{center}

%\vspace{0.7cm}

\vspace{0.5cm}
\noindent\textbf{Palabras clave}: \textit{software libre}
\vspace{0.7cm}

\noindent\textbf{Resumen}\\
	

\cleardoublepage

\begin{center}
	{\large\bfseries Same, but in English}\\
\end{center}
\begin{center}
	Student's name\\
\end{center}
\vspace{0.5cm}
\noindent\textbf{Keywords}: \textit{open source}, \textit{floss}
\vspace{0.7cm}

\noindent\textbf{Abstract}\\


\cleardoublepage

\thispagestyle{empty}

\noindent\rule[-1ex]{\textwidth}{2pt}\\[4.5ex]

D. \textbf{Tutora/e(s)}, Profesor(a) del ...

\vspace{0.5cm}

\textbf{Informo:}

\vspace{0.5cm}

Que el presente trabajo, titulado \textit{\textbf{Chief}},
ha sido realizado bajo mi supervisión por \textbf{Estudiante}, y autorizo la defensa de dicho trabajo ante el tribunal
que corresponda.

\vspace{0.5cm}

Y para que conste, expiden y firman el presente informe en Granada a Junio de 2018.

\vspace{1cm}

\textbf{El/la director(a)/es: }

\vspace{5cm}

\noindent \textbf{(nombre completo tutor/a/es)}

\chapter*{Agradecimientos}

Poner aquí agradecimientos...


%\frontmatter
%\tableofcontents
%\listoffigures
%\listoftables
%
%\mainmatter
%\setlength{\parskip}{5pt}

\chapter{Introducción}
En esta sección se detalla la importancia de la gestión económica personal y cómo el avance de la digitalización ha transformado la manera en que interactuamos con nuestras finanzas. Se presenta un recorrido por la evolución de los servicios financieros digitales en España, desde los primeros cajeros automáticos hasta las modernas formas de pago como \textit{Bizum} y \textit{NFC}. Asimismo, se analizan los retos actuales que enfrentan los usuarios al gestionar sus finanzas personales, dadas las múltiples fuentes de ingresos y métodos de pago disponibles. Se explora cómo las soluciones \textit{Fintech} han surgido para abordar estos desafíos.

\section{Contexto. Descripción del problema} 
Según su definición en la \href{https://dle.rae.es/economía}{RAE} la economía es la 
\textit{Administración eficaz y razonable de los bienes}, dicho esto y teniendo en 
cuenta que nos concierne a todos, es crucial tener control sobre ella. 
De manera individual, las decisiones financieras influyen en la calidad de vida de las personas,
un uso adecuado es esencial para evitar problemas como el endeudamiento, la 
falta de ahorro o la incapacidad de cumplir metas económicas.

La transición hacia lo digital está cambiando nuestra forma de interactuar con 
la información, nos permite realizar tareas de forma remota, rápida y eficiente. 
En los últimos años, en España la banca ha evolucionado rápidamente hacia lo digital, comenzando en los años 70 con la introducción de cajeros automáticos. En los 80, las tecnologías de la información mejoraron la eficiencia de los mercados financieros, y en los 90, surgieron los primeros servicios bancarios a distancia, como la banca telefónica. En 1995, se lanzó el primer software que permitía ver finanzas online, y en 1999 los bancos españoles comenzaron a ofrecer sus primeros servicios de consulta online. A partir del año 2006 la llegada de los smartphones y la evolución de Internet aceleraron la transformación \cite{hebrero2022fintech}.

Actualmente, esta digitalización ha brindado la posibilidad de utilizar diversas formas de de pago como transferencias bancarias, el uso de tarjetas de crédito, sistemas de pago instantáneo como \href{https://bizum.es/}{Bizum}, pago con dispositivos como smartphones y relojes inteligentes mediante NFC \footnote{NFC (Near Field Communication) es una tecnología que permite una comunicación de corto alcance entre dispositivos inhalámbricos para intercambiar pequeñas cantidades de datos} y un largo etcétera. Por supuesto, no podemos olvidar el uso de dinero en efectivo, que sigue siendo una forma de pago común. 

En una época en la que los precios son muy elevados, pero a la vez el consumo parece estar desvocado, es importante administrar el dinero de manera responsable. Esto implica planificación y control sobre cómo y dónde lo gastamos.
Sin embargo, no es fácil hacer el seguimiento de las finanzas personales ya que se 
juntan diversos factores a tener en cuenta: podemos tener varias fuentes de ingresos y/o 
gastos, planes de ahorro, diferentes cuentas bancarias, uso de dinero tanto 
efectivo como digital, etc. por lo que en ocasiones resulta tedioso 
llevar las cuentas al día para analizar correctamente nuestro consumo.

Como ayuda para solventar este problema, han aparecido aplicaciones que nos ayudan a gestionar nuestros gastos, ahorros, inversiones, etc. pero como ocurre en cualquier ámbito, no siempre se adaptan a las necesidades concretas de algunos usuarios. Estas herramientas que usan la tecnología para ofrecer a los usuarios servicios financieros, forman parte de las denominadas \textit{Fintech} \cite{schueffel2016taming}.

\section{Motivación}
Dado el escenario descrito, la motivación para realizar este proyecto surge del deseo de contribuir a la mejora en la calidad de vida de las personas a través de un mayor control sobre sus decisiones financieras. 

Ante la diversidad de métodos de pago mencionada y la necesidad de una gestión eficiente de las finanzas, este proyecto propone unificar la información de todas las fuentes de ingresos y gastos del usuario. Al ofrecer una visión global y realista de su situación económica, la aplicación permitirá a los usuarios tener un control centralizado sobre sus finanzas, facilitando la toma de decisiones informadas y mejorando su capacidad para planificar, anticiparse y gestionar sus recursos financieros de manera más efectiva.

Se centra en el control del dinero analizando el consumo mes a mes y en el ahorro guiado por la planificación, anticipación y previsión de gastos. Para ello 
se introduce el concepto de un monedero virtual, que cada mes empieza de cero, se va llenando con los ingresos y se vacía con los gastos y los ahorros dedicados a otros objetivos.\\

Este proyecto es software libre, y está liberado con la licencia \cite{gplv3}.\\
Puede encontrarse en el siguiente repositorio:
\href{https://github.com/elenaortegacontreras/TFG}{elenaortegacontreras/TFG}, 
donde se ha desarrollado en abierto desde su inicio.

\section{Objetivos} \label{sect:goals}
\textit{El objetivo principal de este proyecto es desarrollar una aplicación web (\textbf{SIGMA}: Sistema Inteligente de Gestión Monetaria Automatizada) 
para la gestión financiera. Usando tecnologías de reconocimiento
óptico de caracteres (OCR) y geolocalización, permite al usuario llevar a cabo un 
análisis de gastos detallado y facilita una gestión eficiente de su dinero
a través de un monedero virtual, introduciendo los datos de 
forma sencilla.}

Para el cumplimiento de este objetivo general se plantean los siguientes objetivos específicos:
\begin{enumerate}
    \item Analizar las herramientas de gestión financiera existentes en el mercado para su adecuada aplicación en el desarrollo del proyecto.  
    \item Plantear un diseño escalable y desarrollar una arquitectura de software modular para facilitar la integración de futuras funcionalidades o mejoras.
    \item Diseñar e implementar la interfaz de usuario para permitir la creación 
         de transacciones (ingresos, gastos y ahorros), categorías de gasto y objetivos de ahorro suficientemente personalizables. \label{obj:O3}
    \item Implementar herramientas de visualización de datos por medio de mapas, gráficos y resúmenes automatizados para analizar los patrones de gasto y conseguir un seguimiento claro de los mismos por parte del usuario.\label{obj:O4}
    \item Estudiar e integrar tecnología de reconocimiento óptico de caracteres (OCR) y de reconocimiento de patrones en un texto para escanear tickets de compra; como punto de partida al procesamiento del texto obtenido. \label{obj:05}
    \item Implementar un sistema que permita automatizar la inserción de gastos en la aplicación, extrayendo datos relevantes de los tickets de compra para que la interverción por parte del usuario en el proceso sea mínima.\label{obj:06}
    \item Facilitar la búsqueda de comercios. Integrar en la aplicación capacidades de geolocalización, para que el usuario pueda identificar en la aplicación el comercio donde ha realizado un gasto suponiéndole el mínimo esfuerzo posible.\label{obj:O7}
    \item Implementar una solución económica y sostenible, para permitir el acceso a la aplicación a un mayor número de usuarios y facilitar su mantenimiento a largo plazo.\label{obj:O8}
\end{enumerate}
\chapter{Estado del arte}
\section{Crítica al estado del arte}

\section{Análisis de las posibles soluciones al problema propuesto}
En la búsqueda de soluciones al problema encontramos varias alternativas. A parte 
de la propia gestión manual, ya sea por la vía tradicional \textit{usando \textbf{papel} y boli} , 
o bien algo más automatizado? como el uso de \textbf{hojas de cálculo}, existen aplicaciones 
que nos facilitan esa tarea e incorporan resúmenes detallados. 

1. En la actualidad la mayoría de \textbf{aplicaciones bancarias} incorporan análisis de gastos. 
Aunque nos permiten ver un resumen mensual de los gastos o incluso por categorías, en 
general no se pueden personalizar demasiado. Algunos métodos de ahorro frecuentes son acumular 
de manera automática los céntimos de euro que sobran al redondear cada compra en una cuenta 
de ahorro virtual, o apartar una cantidad al final del mes. En este caso, nos centraremos 
en el ahorro guiado por la planificación, anticipación y previsión de gastos.

???????????????????
¿Cómo indico que he preguntado a clientes potenciales para la app?
?????????????????

Primero analizamos las características más interesantes (en relación al proyecto) 
de las aplicaciones de los primeros en la lista 
de bancos españoles más grandes https://es.wikipedia.org/wiki/Anexo:Bancos_de_Espa%C3%B1a :

- \textbf{CaixaBank} En la aplicación de \textbf{ImaginBank}, con el servicio MyMonz https://www.imagin.com/app/mymonz 
se recibe un informe mensual de gastos e ingresos que podemos ver por categorías. 
En cuanto a los planes de ahorro, con Mi Hucha podemos crear retos (máximo 5 huchas diferentes) 
https://www.imagin.com/ahorro/retos-ahorro 
con aportaciones periódicas mes a mes o puntuales, pero de forma global.
No aparece la opción de establecer un plan de consumo por categoría del gasto, 
con el que nos adelantemos y seamos conscientes en cada pago del cumplimiento 
de nuestros objetivos. 

- \textbf{Banco Santander}
https://www.bancosantander.es/particulares/banca-online/apps/santander
La aplicación de este banco tiene una zona de "Análisis de gastos", que nos crea 
informes similares a los de ImaginBank. 
Por otro lado, ésta sí nos permite establecer un 
presupuesto por categorías que podemos revisar eventualmente para guiarnos en el 
cumplimiento de nuestro objetivo y recibir alertas al acercarnos al límite.

- \textbf{BBVA}
Esta última parece ser la más completa en cuanto a opciones de ahorro y análisis de gastos.

Con "Mi día a día" https://www.bbva.es/personas/banca-online/control-gastos-mi-dia-a-dia.html 
obtenemos, al igual que en el resto de aplicaciones, el resumen con los movimientos de la cuenta.
Respecto a los métodos para ahorrar existe una cuenta gratuita denominada 
cuenta metas https://www.bbva.es/personas/productos/cuentas/cuenta-ahorro-metas.html#establece-tus-datos-de-contacto-y-acceso 
en la que apartar dinero a modo de hucha para conseguir hasta un máximo de 5 objetivos. 

El apartado presupuestos https://www.bbva.es/general/salud-financiera/economia-domestica/gestor-de-gastos-y-presupuestos.html 
es el lugar donde definir una cantidad máxima de gasto deseada 
al mes y por categoría, mide con códigos de colores la evolución del consumo 
respecto al presupuesto.
Encontramos también Apartados https://www.bbva.es/finanzas-vistazo/tu-guia-bbva/app/apartados-una-nueva-forma-de-ahorrar.html#:~:text=Apartados%2C%20de%20BBVA%2C%20es%20un,manera%20m%C3%A1s%20c%C3%B3moda%20y%20eficiente. 
donde el objetivo es el mismo que en presupuestos, con la diferencia de 
que es un espacio para separar visualmente tu dinero según tus necesidades,
apartando la cantidad que quieras para cada tipo de gasto.

2. Además, los gastos y configuración que se hacen en la aplicación que ofrece 
el banco, generalmente no se pueden exportar a otras aplicaciones,
por lo que si se quiere cambiar de banco, se pierde toda la información, 
y con ella los planes de ahorro.

3. Por otro lado, existen aplicaciones de terceros que permiten unificar los gastos 
independientemente del banco al que pertenezca el usuario. Entre las más conocidas 
encontramos:
---------- cosas que investigar de cada una de las aplicaciones ----------
- coste 
- cómo se introducen los datos
- necesita credenciales de acceso al banco?
- hace analisis?
- exportación de datos??
- permite hacer objetivos por categorías?

--> 1a web caracteristicas destacables: https://n26.com/es-es/blog/9-apps-para-controlar-tus-gastos 
--> 2da web: https://www.20minutos.es/tecnologia/aplicaciones/8-apps-para-gestionar-tu-economia-y-ahorrar-a-final-de-mes-4939558/ 
--> 3a web: https://www.bbva.com/es/salud-financiera/las-10-apps-para-gestionar-y-compartir-tus-gastos/ 
--> 2da web: https://www.ionos.es/digitalguide/online-marketing/vender-en-internet/las-mejores-apps-para-controlar-tus-gastos/ 

\textbf{Fintonic} https://www.fintonic.com/es-ES/inicio/ 
Un aspecto que puede causar reticencia a usar la aplicación es que necesita 
las claves de acceso a la cuenta bancaria (aunque sean de solo lectura) 
para acceder a los datos relacionados con los movimientos de la cuenta.
Tiene una interfaz intuitiva, y recoge, clasifica y analiza los gastos.


\textbf{Money Manager}  https://moneymanagerapp.com/
Se deben introducir los datos manualmente, aunque también se pueden importar
Solo se encuentra disponible en inglés. Ofrece una gran cantidad de opciones
para analizar gastos e incluye la opción de añadir fotos a las transacciones.

\textbf{Moneyfy} 

\textbf{Wallet} 

\textbf{1Money} 
Muy completa, pero solo disponible para android.


\textbf{Money Hero} https://moneyhero.site/es  


siguiente tarea:
AHORA ANALIZAR PARA CADA OPCIÓN LAS APLICACIONES MÁS DESTACABLLES E INCONVENIENTES.


\textbf{Solución planteada / LO QUE VA A BUSCAR MI APP:}
Para guiarme en lo que quiero que tenga y por qué:
https://www.bbva.es/finanzas-vistazo/ef/ahorro/como-hacer-para-ahorrar-dinero-y-no-gastarlo.html

Ninguna de las opciones invetigadas previamente incluye todos las funcionalidades 
que se proponen como requisitos en este proyecto, por lo que se plantea la creación
de una aplicación que englobe todas ellas.

opción de establecer un plan de consumo por categoría del gasto, con el que nos adelantemos y 
seamos conscientes en cada pago del cumplimiento de nuestros objetivos.a

Proponiendo las siguientes soluciones a cada uno de los problemas descritos:

Siendo la meta mejorar la salud financiera de los usuarios, se plantea una aplicación que:

Respondiendo a cada inconveniente antes descrito se plantea una aplicación que 
agrupe las siguientes soluciones / características?:


1 --> ADAPTABILIDAD Y PERSONALIZACIÓN: 
- Lo de establecer planes de ahorro por categoría decir que lo hace
Adaptándose así a las necesidades de cada usuario.

2 -3 --> UNICIDAD / UNIFICACIÓN?
Con objetivo de solucionar la pérdida de la información al cambiar de banco, 
se propone optar' por una \textbf{aplicación externa al banco}. Para la unificación de gastos 
con tarjeta, pagos en efectivo, transferencias y demás operaciones, la 
aplicación propuesta permitirá añadir gastos de forma manual, y mediante 
escaneo de tickets. 



4 --> AUTOENGAÑO
Monedero virtual que almacena parte de los ingresos.


\textbf{ANÁLISIS DE LAS APLICACIONES ACTUALES}
\textbf{Aplicaciones en el mercado que resuelven el problema}
Entre las aplicaciones más conocidas para la gestión de gastos se encuentran:
PLEO, REVOLUT, N26, Bnext, Fintonic, Money Pro, Spendee, Wallet, Money Manager, Money Lover, etc.

\textbf{Pleo} De pago. Prueba gratuita.

\textbf{Revolut} De pago. Prueba gratuita.

\textbf{Fintonic} Gratuita. 

\textbf{Problemas de las aplicaciones actuales}







\textbf{ELECCIÓN DE HERRAMIENTAS PARA EL DESARROLLO}
\textbf{Tecnología para el desarrollo de la aplicación}
Para desarrollar una aplicación multiplataforma, se pueden usar tecnologías como Flutter, React Native, Xamarin, etc.
https://flutter.dev/ :"Flutter es un marco de código abierto de Google para crear aplicaciones multiplataforma (móviles, 
web, de escritorio) de forma nativa a partir de una única base de código".


OCR:
necesitamos BD, reconocimiento de espaciado entre datos
por escaneo es lo más fiable, luego viene el uso de la cámara
Aclarar la imagen, lo mejor es una imagen en blanco y negro para que sea lo más sencillo

Google tiene un servicio para extraer info de un PDF no editable con DocumentAI
Hay librerías en python para usar OCR:
    - pytesseract, para reconocer caracteres 

- easyOCR
- keras-OCR
- trOCR
- docTR

\section{Propuesta}


%
%\input{capitulos/02_EspecificacionRequisitos}
%
%\input{capitulos/03_Planificacion}
%
%\input{capitulos/04_Analisis}
%
%\input{capitulos/05_Diseno}
%
%\input{capitulos/06_Implementacion}
%
%\input{capitulos/07_Pruebas}
%
%\input{capitulos/08_Conclusiones}
%
%%\chapter{Conclusiones y Trabajos Futuros}
%
%
%%\nocite{*}
%\bibliography{bibliografia/bibliografia}\addcontentsline{toc}{chapter}{Bibliografía}
%\bibliographystyle{miunsrturl}
%
\bibliography{bibliografia}\addcontentsline{toc}{chapter}{Bibliografía}
%\appendix
%\input{apendices/manual_usuario/manual_usuario}
%%\input{apendices/paper/paper}
%\input{glosario/entradas_glosario}
% \addcontentsline{toc}{chapter}{Glosario}
% \printglossary
\chapter*{}
\thispagestyle{empty}

\end{document}
